\begin{register}{H}{AES\_KEY\_\regindex{n}\_REG (\regindex{n}: 0-7)}{0x0000+4*\regindex{n}}\label{regdesc:AESKEYnREG}
\regfield{AES\_KEY\_\regindex{n}\_REG (\regindex{n}: 0-7)}{32}{0}{{0x000000000}}%
\reglabel{Reset}\regnewline%
\begin{regdesc}\begin{reglist}
\item [AES\_KEY\_\regindex{n}\_REG (\regindex{n}: 0-7)] 表示 AES 密钥资料。(R/W)
\end{reglist}\end{regdesc}
\end{register}

\begin{register}{H}{AES\_TEXT\_IN\_\regindex{m}\_REG (\regindex{m}: 0-3)}{0x0020+4*\regindex{m}}\label{regdesc:AESTEXTINnREG}
\regfield{AES\_TEXT\_IN\_\regindex{m}\_REG (\regindex{m}: 0-3)}{32}{0}{{0x000000000}}%
\reglabel{Reset}\regnewline%
\begin{regdesc}\begin{reglist}
\item [AES\_TEXT\_IN\_\regindex{m}\_REG (\regindex{m}: 0-3)] 表示 Typical AES 工作模式下的输入文本。(R/W)
\end{reglist}\end{regdesc}
\end{register}


\begin{register}{H}{AES\_TEXT\_OUT\_\regindex{m}\_REG (\regindex{m}: 0-3)}{0x0030+4*\regindex{m}}\label{regdesc:AESTEXTOUTnREG}
\regfield{AES\_TEXT\_OUT\_\regindex{m}\_REG (\regindex{m}: 0-3)}{32}{0}{{0x000000000}}%
\reglabel{Reset}\regnewline%
\begin{regdesc}\begin{reglist}
\item [AES\_TEXT\_OUT\_\regindex{m}\_REG (\regindex{m}: 0-3)] 表示 Typical AES 工作模式下的输出文本。(RO)
\end{reglist}\end{regdesc}
\end{register}


\begin{register}{H}{AES\_MODE\_REG}{0x0040}\label{regdesc:AESMODEREG}
\regfield{(reserved)}{29}{3}{{0x00000000}}%
\regfield{AES\_MODE}{3}{0}{{0}}%
\reglabel{Reset}\regnewline%
\begin{regdesc}\begin{reglist}
\item [AES\_MODE] 配置 AES 加速器的密钥长度和加解密方向。\\
0：AES-128 加密\\
1：保留\\
2：AES-256 加密\\
3：保留\\
4：AES-128 解密\\
5：保留\\
6：AES-256 解密\\
7：保留\\
(R/W)
\end{reglist}\end{regdesc}
\end{register}


\begin{register}{H}{AES\_DMA\_ENABLE\_REG}{0x0090}\label{regdesc:AESDMAENABLEREG}
\regfield{(reserved)}{31}{1}{{0x00000000}}%
\regfield{AES\_DMA\_ENABLE}{1}{0}{{0}}%
\reglabel{Reset}\regnewline%
\begin{regdesc}\begin{reglist}
\item [AES\_DMA\_ENABLE] 配置 AES 加速器的工作模式。\\
0：Typical AES\\
1：DMA-AES\\
(R/W)
\end{reglist}\end{regdesc}
\end{register}


\begin{register}{H}{AES\_BLOCK\_MODE\_REG}{0x0094}\label{regdesc:AESBLOCKMODEREG}
\regfield{(reserved)}{29}{3}{{0x00000000}}%
\regfield{AES\_BLOCK\_MODE}{3}{0}{{0}}%
\reglabel{Reset}\regnewline%
\begin{regdesc}\begin{reglist}
\item [AES\_BLOCK\_MODE] 配置 AES 加速器在 DMA-AES 工作模式下的块模式。\\
0：ECB (Electronic Code Block)\\
1：CBC (Cipher Block Chaining)\\
2：OFB (Output FeedBack)\\
3：CTR (Counter)\\
4：CFB8 (8-bit Cipher FeedBack)\\
5：CFB128 (128-bit Cipher FeedBack)\\
6：保留\\
7：保留\\
(R/W)
\end{reglist}\end{regdesc}
\end{register}


\begin{register}{H}{AES\_BLOCK\_NUM\_REG}{0x0098}\label{regdesc:AESBLOCKNUMREG}
\regfield{AES\_BLOCK\_NUM}{32}{0}{{0x00000000}}%
\reglabel{Reset}\regnewline%
\begin{regdesc}\begin{reglist}
\item [AES\_BLOCK\_NUM] 代表 DMA-AES 工作模式下待加解密的文本块数。详情请见章节 \ref{sec:aes-block-number}。(R/W)
\end{reglist}\end{regdesc}
\end{register}


\begin{register}{H}{AES\_INC\_SEL\_REG}{0x009C}\label{regdesc:AESINCSELREG}
\regfield{(reserved)}{31}{1}{{0x00000000}}%
\regfield{AES\_INC\_SEL}{1}{0}{{0}}%
\reglabel{Reset}\regnewline%
\begin{regdesc}\begin{reglist}
\item [AES\_INC\_SEL] 配置 CTR 块模式使用的标准增量函数。\\
0：INC$_{32}$ 标准增量函数\\
1：INC$_{128}$ 标准增量函数\\
(R/W)
\end{reglist}\end{regdesc}
\end{register}

\begin{register}{H}{AES\_TRIGGER\_REG}{0x0048}\label{regdesc:AESTRIGGERREG}
\regfield{(reserved)}{31}{1}{{0x00000000}}%
\regfield{AES\_TRIGGER}{1}{0}{x}%
\reglabel{Reset}\regnewline%
\begin{regdesc}\begin{reglist}
\item [AES\_TRIGGER] 配置是否启动 AES 运算。\\
0：无效果\\
1：启动\\
(WT)
\end{reglist}\end{regdesc}
\end{register}


\begin{register}{H}{AES\_STATE\_REG}{0x004C}\label{regdesc:AESSTATEREG}
\regfield{(reserved)}{30}{2}{{0x00000000}}%
\regfield{AES\_STATE}{2}{0}{{0x0}}%
\reglabel{Reset}\regnewline%
\begin{regdesc}\begin{reglist}
\item [AES\_STATE] 代表 AES 加速器的状态。\\
在 Typical AES 工作模式下，\\
0：IDLE \\
1：WORK\\
2：无效果\\
3：无效果\\
在 DMA-AES 工作模式下，\\
0：IDLE\\
1：WORK\\
2：DONE\\
3：无效果\\
(RO) 
\end{reglist}\end{regdesc}
\end{register}

\begin{register}{H}{AES\_DMA\_EXIT\_REG}{0x00B8}\label{regdesc:AESDMAEXITREG}
\regfield{(reserved)}{31}{1}{{0x00000000}}%
\regfield{AES\_DMA\_EXIT}{1}{0}{x}%
\reglabel{Reset}\regnewline%
\begin{regdesc}\begin{reglist}
\item [AES\_DMA\_EXIT] 配置是否退出 AES 运算。\\
0：无效果\\
1：退出\\
仅在 DMA-AES 模式下有效。(WO)
\end{reglist}\end{regdesc}
\end{register}


\begin{register}{H}{AES\_PSEUDO\_REG}{0x00D0}\label{regdesc:AESPSEUDOREG}
\regfield{(reserved)}{22}{10}{0000000000000000000000}%
\regfield{AES\_PSEUDO\_RNG\_CNT}{3}{7}{{7}}%
\regfield{AES\_PSEUDO\_INC}{2}{5}{{2}}%
\regfield{AES\_PSEUDO\_BASE}{4}{1}{{2}}%
\regfield{AES\_PSEUDO\_EN}{1}{0}{{0}}%
\reglabel{Reset}\regnewline%
\begin{regdesc}\begin{reglist}
\label{fielddesc:AESPSEUDOEN}\item [AES\_PSEUDO\_EN] 配置是否启用 AES 伪轮函数。\\
0：禁用\\
1：启用\\
(R/W)
\label{fielddesc:AESPSEUDOBASE}\item [AES\_PSEUDO\_BASE] 配置伪轮次的基本次数。(R/W)
\label{fielddesc:AESPSEUDOINC}\item [AES\_PSEUDO\_INC] 配置伪轮次的随机增量次数。(R/W)
\label{fielddesc:AESPSEUDORNGCNT}\item [AES\_PSEUDO\_RNG\_CNT] 配置伪轮函数中随机密钥更新的频率。(R/W)
\end{reglist}\end{regdesc}
\end{register}


\begin{register}{H}{AES\_INT\_CLR\_REG}{0x00AC}\label{regdesc:AESINTCLRREG}
\regfield{(reserved)}{31}{1}{{0x00000000}}%
\regfield{AES\_INT\_CLR}{1}{0}{x}%
\reglabel{Reset}\regnewline%
\begin{regdesc}\begin{reglist}
\item [AES\_INT\_CLR] 配置是否清除 AES 中断。\\
0：无效果\\
1：清除\\
(WT)
\end{reglist}\end{regdesc}
\end{register}


\begin{register}{H}{AES\_INT\_ENA\_REG}{0x00B0}\label{regdesc:AESINTENAREG}
\regfield{(reserved)}{31}{1}{{0x00000000}}%
\regfield{AES\_INT\_ENA}{1}{0}{0}%
\reglabel{Reset}\regnewline%
\begin{regdesc}\begin{reglist}
\item [AES\_INT\_ENA] 配置是否使能 AES 中断功能。\\
0：禁用\\
1：使能\\
(R/W)
\end{reglist}\end{regdesc}
\end{register}
